% META: {"id": "1NSI-AGT-EVAL-A", "chapitre": "1NSI-ALGO-PARCOURS-TRIS", "type_objet": "evaluation", "capacites": ["P-ALGO-01A", "P-ALGO-01B", "P-ALGO-02A", "P-ALGO-02C"], "mode_creation": "ex_nihilo", "sources_inspiration": [], "duree_min": 45, "status": "verified"}
\section*{Évaluation A --- Parcours et tris}
\textit{Durée : 45 minutes. Ordinateur avec interpréteur Python autorisé.}

\baremeIndicatif{Ex. 1 : 10 pts (recherche, extremum, moyenne) ; Ex. 2 : 10 pts (tri par insertion, invariant)}

\bigskip

\textbf{Exercice 1} --- \textit{Parcours séquentiel} \hfill (10 points)

On donne le tableau \lstinline{[22, 15, 8, 30, 11]}.

\begin{enumerate}
  \item (3 pts) À quel indice se trouve la valeur $30$ ? Utiliser \lstinline{recherche_occurrence} du cours.
  \item (3 pts) Calculer le maximum de ce tableau avec \lstinline{maximum}.
  \item (4 pts) Calculer la moyenne de ce tableau avec \lstinline{moyenne}. Détailler le calcul à la main.
\end{enumerate}

% BEGIN-VERIFY
% def recherche_occurrence(tableau, valeur):
%     for i in range(len(tableau)):
%         if tableau[i] == valeur:
%             return i
%     return -1
% def maximum(tableau):
%     maxi = tableau[0]
%     for x in tableau[1:]:
%         if x > maxi:
%             maxi = x
%     return maxi
% def moyenne(tableau):
%     return sum(tableau) / len(tableau)
% t = [22, 15, 8, 30, 11]
% assert recherche_occurrence(t, 30) == 3
% assert maximum(t) == 30
% assert moyenne(t) == 17.2
% END-VERIFY

\bigskip

\textbf{Exercice 2} --- \textit{Tri par insertion} \hfill (10 points)

On applique le tri par insertion au tableau \lstinline{[9, 2, 6, 4]}.

\begin{enumerate}
  \item (5 pts) Dérouler l'algorithme : donner le contenu du tableau après chaque itération de la boucle externe.
  \item (3 pts) Vérifier, à chaque étape, l'invariant de boucle du cours (le préfixe traité est trié).
  \item (2 pts) Quel est le coût de cet algorithme dans le pire cas, en fonction de la taille $n$ du tableau ?
\end{enumerate}

% BEGIN-VERIFY
% def tri_insertion(tableau):
%     tableau = tableau[:]
%     for i in range(1, len(tableau)):
%         valeur = tableau[i]
%         j = i - 1
%         while j >= 0 and tableau[j] > valeur:
%             tableau[j + 1] = tableau[j]
%             j -= 1
%         tableau[j + 1] = valeur
%     return tableau
% assert tri_insertion([9, 2, 6, 4]) == [2, 4, 6, 9]
% END-VERIFY
